% google-fonts.tex
\usepackage{fontspec}
\setmainfont{Roboto}

\newif\ifxetexorluatex
\ifx\XeTeXrevision\undefined
  \ifx\directlua\undefined
    \xetexorluatexfalse
  \else
    \xetexorluatextrue  % LuaLaTeX
  \fi
\else
  \xetexorluatextrue    % XeLaTeX
\fi

\ifxetexorluatex
  % --- XeLaTeX / LuaLaTeX ---
  \usepackage{fontspec}
  % Dossier où tu mets les TTF/OTF téléchargées sur fonts.google.com
  % Exemple avec Roboto + Roboto Mono
  \newfontfamily\Roboto{Roboto}[
    Path=fonts/Roboto/, Extension=.ttf,
    UprightFont=*-Regular, BoldFont=*-Bold,
    ItalicFont=*-Italic, BoldItalicFont=*-BoldItalic
  ]
  \newfontfamily\RobotoMono{RobotoMono}[
    Path=fonts/RobotoMono/, Extension=.ttf,
    UprightFont=*-Regular, BoldFont=*-Bold,
    ItalicFont=*-Italic, BoldItalicFont=*-BoldItalic
  ]
  \setmainfont{Roboto}      % texte principal
  \setsansfont{Roboto}      % sans-serif
  \setmonofont{RobotoMono}  % typewriter/code
\else
  % --- pdfLaTeX (fallback, pas de Google Fonts) ---
  \usepackage[T1]{fontenc}
  \usepackage[utf8]{inputenc}
  \usepackage{lmodern}      % ou newtxtext,newtxmath si tu préfères
\fi
